% 二、问题分析
\section{问题分析}
\subsection{总体分析：信息状态与库存机会成本}
四个问题具有相同的物理约束。每一时段均需满足外网购电、光伏供能和储能放电与负载、储能充电之间的能量平衡；储电量按照充放电动作递推，并受容量与功率上限约束。随着题目推进，变化主要体现在决策前可获得的信息逐步丰富，以及购电计划由 0:00 一次确定扩展为允许日内调整。

储能状态使各时段决策相互关联。当前释放的库存能够降低本时段购电需求，同时会减少未来高成本缺口时可利用的电量。因此，仅记录一条计划库存轨迹不足以支持随机环境下的实时执行，还需要刻画不同库存状态对应的未来最低费用。

记 $\mathcal I_t$ 为时刻 $t$ 的信息状态，其中包含截至当前已经获得的负载、光伏和价格观测，已经发布的预测信息，当前储电量以及当前允许使用的交易权限。在给定 $\mathcal I_t$ 的条件下，以库存 $e$ 为状态构造未来费用函数，由其边际变化衡量库存电量的跨时段机会成本。问题 1 得到完整信息下的确定性价值函数 $F_t(e)$；问题 2 将其扩展为预测场景下的平均未来费用 $\bar H_t(e)$；问题 3 进一步随预报发布时刻 $\tau$ 更新；问题 4 再纳入未来价格信息。由此，四问可以在同一能量调度结构下随信息状态逐步扩展。

\subsection{问题 1 的分析}
问题 1 给定一天内用于决策的电价、负载与光伏预测曲线，各时段购电量、充放电量和库存之间均为线性关系，目标函数为各时段电价与购电量之和，因此建立确定性线性规划，在能量平衡、充放电功率、库存边界及日初日末库存相等的约束下求得全天计划的最低购电费用。

线性规划能够给出一组全局最优调度结果，但后三问还要处理执行过程中库存偏离计划值后的重新决策。为此，以当前储电量 $e$ 为状态，定义从时段 $t$ 开始至计划结束的剩余最低费用 $F_t(e)$，并建立连续状态 Bellman 递推。由于单段动作成本具有凸分段线性结构，$F_t(e)$ 在有效域上保持凸分段线性，可用折线断点和斜率精确表示，无需对储电量进行网格离散。LP 承担确定全天计划调度和最优费用的作用，连续状态 DP 提供库存状态的未来价值表示；两种独立求解路径的最优值一致性又构成对优化核的交叉验证。

\subsection{问题 2 的分析}
问题 2 中负载与光伏出力在 0:00 制定计划时尚未完全实现，需要同时处理计划采购的不确定性和实际运行阶段的储能控制。预测层利用附件 2 的历史数据描述可预测结构：将负载拆解为日电量水平与归一化日内形状，分别预测后重构全天计划曲线；光伏则使用决策时点以前已经完成的历史日构造基准预测。

预测误差在一天内具有连续性，负载与光伏误差之间也可能存在共同变化，因此场景层以历史完整日的预测残差为基本单元，并保留同一历史日的负载--光伏残差配对关系。决策时刻选取最近 $M=30$ 条已经完整实现的历史残差路径，叠加到当前点预测后形成经验场景集。

0:00 计划量一经制定便在全天计划内冻结，具有"先确定采购、后观察随机实现"的两阶段结构。采购层采用样本平均近似\cite{saa-book,kleywegt}，在多个历史场景下共同优化第一阶段计划购电量；场景内第二阶段变量用于评价储能追索和潜在紧急购电。由于该代理问题允许场景路径完整可见，其最优值用于计划选择与机会成本估计，不直接作为真实执行动作。

真实执行只使用当前已经观测的信息。对固定计划和每条场景逆向计算未来费用，并取等权平均得到 $\bar H_t(e)$。当出现供能缺口时，将当段紧急购电边际成本与未来库存边际价值进行比较，得到动态库存保留水平 $R_t$：库存高于该水平的可用部分可以释放，剩余部分为未来潜在高成本缺口保留。为减弱 24:00 截断造成的终端效应，主方法将前瞻范围扩展到次日，并以次日自由购电的价值函数给午夜库存定价。全年实际费用最终由真实负载、真实光伏和真实结算规则重新计算。

\subsection{问题 3 的分析}
问题 3 同时增加日内预报信息和购电调整权限。0:00 制定的原计划作为全天基准；6:00、12:00 和 18:00 获得新的未来 24\,h 光伏预报后，剩余净负荷分布和库存机会成本需要随最新信息重新评估。

在每个预报发布时刻，以真实当前储电量作为新的初始状态，利用截至该时刻已经发布的光伏预报和历史误差场景重新优化尚未执行的购电量，并在目标函数中统一计入调增、调减、紧急购电和跨日库存价值。本轮只正式提交至下一次预报更新时间之前的购电决策，更远时段保留为前瞻计划，待下一次预报到来后重新计算，从而保持整个滚动过程的时序一致性。

为回答是否需要引入其他时刻预报，固定预测方法、场景构造、终端价值和结算规则，按照 $\{0\}$、$\{0,6\}$、$\{0,6,12\}$、$\{0,6,12,18\}$ 的顺序逐步开放日内更新时刻，并以全年真实账单比较相邻方案。由于新增一个发布时刻同时带来新的光伏信息和一次对应的购电调整机会，相邻账单差表示该更新时刻在既定开放顺序下的联合边际实际收益。现有结果表明，12:00 更新的收益最高，其原因与午间光伏预报误差明显收缩以及当天仍保留较长的剩余执行窗口共同相关。

\subsection{问题 4 的分析}
问题 4 中未来电价也需要在购电决策前估计，库存未来费用因此同时取决于未来净负荷与价格路径。根据附件 4 的全年价格结构，将每日电价分解为日价格水平和归一化日内形状，前者刻画跨日整体价格变化，后者描述日内相对稳定的峰谷结构。

联合场景构造时，将价格预测残差与负载、光伏残差按照同一历史信息状态组成三通道完整路径，避免独立抽样破坏变量间的经验相关关系。4-2 沿用问题 2 的交易权限，0:00 计划购电量保持冻结，执行阶段随已经观测的价格信息更新库存价值；4-3 沿用问题 3 的交易权限，在各发布时刻同步更新光伏预报、价格预测与剩余采购。

当全部历史与实际价格、紧急及调整交易价格以及续存费率同步乘以同一正常数时，各优化问题的目标函数整体同比缩放而可行域保持不变，因此最优采购解集、库存保留水平和逐段执行动作均保持不变，费用按相同比例变化。这一齐次性质说明，统一改变绝对价格水平不会改变策略形状，真正影响调度时机的是相对价格结构及其可预测性。

\subsection{总体研究框架}
全文在纵向上依次经过预测与场景构造、采购优化、因果执行和实际结算；在横向上随着问题推进，决策环境由确定信息扩展至供需随机性、日内信息更新和价格不确定性。四问共享同一储能物理约束和基本优化结构，库存未来费用函数随信息状态扩展形成：
\begin{equation}
F_t(e)\;\longrightarrow\;\bar H_t(e)\;\longrightarrow\;\bar H_{t\mid\tau}(e)\;\longrightarrow\;\bar H_{t\mid\tau}(e;\,p),
\label{eq:chain}
\end{equation}
其中，$F_t(e)$ 描述问题 1 完整信息条件下的确定性剩余费用；$\bar H_t(e)$ 用于问题 2 的场景平均库存价值；$\bar H_{t\mid\tau}(e)$ 随问题 3 各日内预报时刻更新；$\bar H_{t\mid\tau}(e;\,p)$ 进一步纳入问题 4 的未来价格路径。问题 2--4 采用连续 walk-forward 方式评价，优化模型中的场景代理目标用于制定计划和估计库存机会成本，最终结果统一按真实实现的负载、光伏、电价以及题目规定的交易费用重新结算。

\begin{figure}[!htbp]\centering
\safeincludegraphics[width=0.96\textwidth]{fig_f01_framework}
\caption{四问的预测、采购与执行框架}
\label{fig:framework}
\end{figure}
